% Part: normal-modal-logic
% Chapter: completeness
% Section: modalities-ccs

\documentclass[../../../include/open-logic-section]{subfiles}

\begin{document}

\olfileid[id]{nml}{com}{mod}

\olsection{Modalitas dan Himpunan Lengkap serta Konsisten}

\begin{explain}
Ketika kita membangun suatu model~$\mModel{M^\Sigma}$ yang himpunan dunianya
diberikan oleh himpunan lengkap dan $\Sigma$-konsisten~$\Delta$ dalam suatu
logika modal normal~$\Sigma$, kita juga perlu mendefinisikan relasi
aksesibilitas~$R^\Sigma$ di antara ``dunia-dunia'' semacam itu. Kita ingin agar
relasi aksesibilitas (dan penetapan $V^\Sigma)$ didefinisikan sedemikian rupa
sehingga $\mSat{M^\Sigma}{!A}[\Delta]$ jika dan hanya jika $!A \in
\Delta$. Bagaimana kita harus melakukannya?

Setelah relasi aksesibilitas didefinisikan, definisi kebenaran pada suatu
dunia menjamin bahwa \iftag{prvBox}{$\mSat{M^\Sigma}{\Box!A}[\Delta]$
  jika dan hanya jika $\mSat{M^\Sigma}{!A}[\Delta']$ bagi semua $\Delta'$
  sedemikian sehingga $R^\Sigma\Delta\Delta'$}{$\mSat{M^\Sigma}{\Diamond!A}[\Delta]$
  jika dan hanya jika $\mSat{M^\Sigma}{!A}[\Delta']$ bagi sekurang-kurangnya
  satu $\Delta'$ sedemikian sehingga $R^\Sigma\Delta\Delta'$}. Pembuktian
bahwa $\mSat{M^\Sigma}{!A}[\Delta]$ jika dan hanya jika $!A \in \Delta$
mengharuskan agar hal ini secara khusus benar bagi !!{formula}s yang diawali
operator modal, yakni \iftag{prvBox}{$\mSat{M^\Sigma}{\Box!A}[\Delta]$ jika
  dan hanya jika $\Box !A \in \Delta$}{$\mSat{M^\Sigma}{\Diamond!A}[\Delta]$
  jika dan hanya jika $\Diamond !A \in \Delta$}. Menggabungkan syarat ini
dengan definisi kebenaran pada suatu dunia bagi
\iftag{prvBox}{$\Box !A$}{$\Diamond !A$} menghasilkan:
\[
\iftag{prvBox}{%
  \Box!A \in \Delta \text{ jika dan hanya jika } !A \in \Delta' \text{
    bagi semua $\Delta'$ dengan } R^\Sigma\Delta\Delta'}{%
  \Diamond!A \in \Delta \text{ jika dan hanya jika } !A \in \Delta' \text{ bagi
    sekurang-kurangnya satu } \Delta' \text{ dengan } R^\Sigma\Delta\Delta'}
\]
\iftag{prvBox}{Tinjau arah dari kiri ke kanan: arah ini menyatakan bahwa jika
  $\Box !A \in \Delta$, maka $!A \in \Delta'$ bagi sebarang $!A$ dan sebarang
  $\Delta'$ dengan $R^\Sigma\Delta\Delta'$. Jika kita menetapkan bahwa
  $R^\Sigma\Delta\Delta'$ jika dan hanya jika $!A \in \Delta'$ bagi semua
  $\Box!A \in \Delta$, maka pernyataan tersebut berlaku. Kita dapat menulis
  syarat di sebelah kanan ``jika dan hanya jika'' secara lebih ringkas sebagai:
  $\Setabs{!A}{\Box!A \in \Delta} \subseteq \Delta'$.}{Tinjau arah dari kanan
  ke kiri: arah ini menyatakan bahwa jika $!A \in \Delta'$, maka
  $\Diamond!A \in \Delta$, bagi sebarang $!A$ dan $\Delta'$ dengan
  $R^\Sigma\Delta\Delta'$. Jika kita menetapkan bahwa
  $R^\Sigma\Delta\Delta'$ berlaku setiap kali $\Diamond!A \in \Delta$ bagi
  semua $!A \in \Delta'$, maka pernyataan tersebut berlaku. Kita dapat
  menulis syarat di sebelah kanan ``jika dan hanya jika'' secara lebih ringkas
  sebagai: $\Setabs{\Diamond!A}{!A \in \Delta'} \subseteq \Delta$.}

Jadi, pertanyaannya ialah: apakah definisi $R^\Sigma$ ini benar-benar
menjamin bahwa \iftag{prvBox}{$\Box !A \in \Delta$ jika dan hanya jika
  $\mSat{M^\Sigma}{\Box !A}[\Delta]$? \iftag{prvDiamond}{Apakah definisi itu
    juga menjamin bahwa }{}}{}\iftag{prvDiamond}{$\Diamond !A \in \Delta$
  jika dan hanya jika $\mSat{M^\Sigma}{\Diamond !A}[\Delta]$?}{} Beberapa
hasil berikutnya akan menetapkan hal ini.
\end{explain}

\begin{defn}
  Jika $\Gamma$ merupakan suatu himpunan !!{formula}s, tetapkan
  \begin{align*}
    \Box\Gamma & = \Setabs{\Box !B}{!B \in \Gamma}\\
    \Diamond\Gamma & = \Setabs{\Diamond !B}{!B \in \Gamma}\\
    \intertext{dan}
    \Box^{-1}\Gamma & = \Setabs{!B}{\Box !B \in \Gamma}\\
    \Diamond^{-1}\Gamma & = \Setabs{!B}{\Diamond !B \in \Gamma}\\
  \end{align*}
\end{defn}

Dengan kata lain, $\Box\Gamma$ adalah $\Gamma$ dengan $\Box$ di depan setiap
!!{formula} dalam~$\Gamma$; $\Box^{-1}\Gamma$ terdiri atas semua
!!{formula}s yang diawali~$\Box$ dalam~$\Gamma$, dengan $\Box$ awal tersebut
dihapus. Definisi ini tidak terlalu penting dengan sendirinya, tetapi akan
sangat menyederhanakan notasi.

Perhatikan bahwa $\Box\Box^{-1}\Gamma \subseteq \Gamma$:
\[
\Box\Box^{-1}\Gamma = \Setabs{\Box!B}{\Box!B \in \Gamma}
\]
yakni, himpunan ini tepat merupakan himpunan semua !!{formula}s dalam~$\Gamma$
yang diawali~$\Box$.

\begin{lem}\ollabel{lem:box1}
  Jika $\Gamma \Proves[\Sigma] !A$, maka $\Box\Gamma \Proves[\Sigma] \Box
  !A$.
\end{lem}

\begin{proof}
  Jika $\Gamma \Proves[\Sigma] !A$, maka terdapat $!B_1$, \dots, $!B_k
  \in \Gamma$ sedemikian sehingga $\Sigma \Proves !B_1 \lif (!B_2 \lif \cdots
  (!B_k \lif !A)\cdots)$. Karena $\Sigma$ normal, berdasarkan aturan \RK{},
  $\Sigma \Proves \Box!B_1 \lif (\Box!B_2 \lif \cdots (\Box!B_k \lif
  \Box!A)\cdots)$, dan jelas bahwa $\Box!B_1$, \dots, $\Box!B_k \in
  \Box\Gamma$. Jadi, berdasarkan definisi, $\Box\Gamma \Proves[\Sigma] \Box
  !A$.
\end{proof}

\begin{lem}\ollabel{lem:box2}
  Jika $\Box^{-1} \Gamma \Proves[\Sigma] !A$, maka $\Gamma
  \Proves[\Sigma] \Box!A$.
\end{lem}

\begin{proof}
  Andaikan $\Box^{-1}\Gamma \Proves[\Sigma] !A$; maka berdasarkan
  \olref{lem:box1}, $\Box\Box^{-1}\Gamma \Proves[\Sigma] \Box!A$. Namun, karena
  $\Box\Box^{-1}\Gamma \subseteq \Gamma$, berdasarkan monotonisitas berlaku
  pula $\Gamma \Proves[\Sigma] \Box!A$.
\end{proof}

\iftag{prvBox}{%
% Box is primitive

\begin{prop}\ollabel{prop:box}
  Jika $\Gamma$ lengkap dan $\Sigma$-konsisten, maka $\Box !A \in
  \Gamma$ jika dan hanya jika, bagi setiap himpunan lengkap dan
  $\Sigma$-konsisten~$\Delta$ sedemikian sehingga $\Box^{-1} \Gamma
  \subseteq \Delta$, berlaku $!A \in \Delta$.
\end{prop}

\begin{proof}
  Andaikan $\Gamma$ lengkap dan $\Sigma$-konsisten. Arah ``hanya jika''
  mudah: andaikan $\Box !A \in \Gamma$ dan $\Box^{-1}\Gamma \subseteq
  \Delta$. Karena $\Box!A \in \Gamma$, berlaku $!A \in
  \Box^{-1}\Gamma \subseteq \Delta$, sehingga $!A \in \Delta$.

  Untuk arah ``jika'', kita membuktikan kontraposisinya: andaikan
  $\Box!A \notin \Gamma$. Karena $\Gamma$ lengkap dan $\Sigma$-konsisten,
  himpunan itu tertutup secara deduktif, sehingga $\Gamma
  \Proves/[\Sigma] \Box !A$. Berdasarkan \olref{lem:box2},
  $\Box^{-1}\Gamma \Proves/[\Sigma] !A$. Berdasarkan
  \olref[prf][con]{prop:consistencyfacts}\olref[prf][con]{prop:consistencyfacts-b},
  $\Box^{-1}\Gamma \cup \{ \lnot!A \}$ $\Sigma$-konsisten. Berdasarkan
  Lema Lindenbaum, terdapat suatu himpunan lengkap dan
  $\Sigma$-konsisten~$\Delta$ sedemikian sehingga
  $\Box^{-1}\Gamma \cup \{ \lnot!A \}
  \subseteq \Delta$. Berdasarkan konsistensi, $!A \notin \Delta$.
\end{proof}
}{%
% Box is not primitive, only Diamond is

\begin{prop}\ollabel{prop:diamond}
  Jika $\Gamma$ lengkap dan $\Sigma$-konsisten, maka $\Diamond !A \in
  \Gamma$ jika dan hanya jika terdapat suatu himpunan lengkap dan
  $\Sigma$-konsisten~$\Delta$ sedemikian sehingga $\Diamond\Delta \subseteq
  \Gamma$ dan $!A \in \Delta$.
\end{prop}

\begin{proof}
  Andaikan $\Gamma$ lengkap dan $\Sigma$-konsisten. Bagian dari kanan ke kiri
  mudah: misalkan $\Delta$ lengkap dan $\Sigma$-konsisten sedemikian sehingga
  $\Diamond\Delta \subseteq \Gamma$ dan $!A \in \Delta$. Maka
  $\Diamond!A \in \Diamond\Delta \subseteq \Gamma$, yakni $\Diamond!A
  \in \Gamma$.

  Untuk arah dari kiri ke kanan, andaikan $\Diamond !A \in \Gamma$. Kita
  harus menunjukkan bahwa terdapat suatu himpunan lengkap dan
  $\Sigma$-konsisten~$\Delta$ dengan $!A \in \Delta$ dan
  $\Diamond\Delta \subseteq \Gamma$. Karena $\Diamond !A \in \Gamma$ dan
  $\Gamma$ $\Sigma$-konsisten, $\Box\lnot !A \notin \Gamma$. Karena
  $\Gamma$ lengkap dan $\Sigma$-konsisten, himpunan itu tertutup secara
  deduktif, sehingga $\Gamma \Proves/[\Sigma] \Box\lnot !A$. Berdasarkan
  \olref{lem:box2}, $\Box^{-1}\Gamma \Proves/[\Sigma] \lnot !A$.
  Berdasarkan
  \olref[prf][con]{prop:consistencyfacts}\olref[prf][con]{prop:consistencyfacts-b},
  $\Box^{-1}\Gamma \cup \{ !A \}$ $\Sigma$-konsisten. Berdasarkan Lema
  Lindenbaum, terdapat suatu himpunan lengkap dan
  $\Sigma$-konsisten~$\Delta$ sedemikian sehingga
  $\Box^{-1}\Gamma \cup \{ !A \} \subseteq
  \Delta$. Jelas bahwa $!A \in \Delta$. Untuk melihat bahwa
  $\Diamond\Delta \subseteq \Gamma$, andaikan hal itu tidak berlaku: bagi
  suatu $!B \in \Delta$, berlaku $\Diamond !B \notin \Gamma$. Karena
  $\Gamma$ lengkap dan $\Sigma$-konsisten, maka $\Box\lnot !B \in \Gamma$.
  Akan tetapi, dengan demikian $\lnot !B \in \Box^{-1}\Gamma \subseteq
  \Delta$, dan hal ini akan membuat $\Delta$ tak konsisten.
\end{proof}
}

\begin{lem}\ollabel{lem:box-iff-diamond}
  Andaikan $\Gamma$ dan $\Delta$ lengkap serta $\Sigma$-konsisten. Maka
  $\Box^{-1}\Gamma \subseteq \Delta$ jika dan hanya jika
  $\Diamond\Delta \subseteq \Gamma$.
\end{lem}

\begin{proof}
  Arah ``hanya jika'': andaikan $\Box^{-1}\Gamma \subseteq \Delta$ dan
  $\Diamond!A \in \Diamond\Delta$ (yakni, $!A \in \Delta$). Untuk
  menunjukkan bahwa $\Diamond!A \in \Gamma$, cukup ditunjukkan bahwa
  $\Box\lnot!A \notin \Gamma$, sebab berdasarkan kemaksimalan akan berlaku
  $\lnot\Box\lnot !A \in \Gamma$. Sekarang, jika $\Box\lnot!A \in \Gamma$,
  maka berdasarkan hipotesis $\lnot!A \in \Delta$, bertentangan dengan
  konsistensi $\Delta$ (sebab $!A \in \Delta$). Jadi,
  $\Box\lnot!A \notin \Gamma$, seperti yang diperlukan.

  Arah ``jika'': andaikan $\Diamond\Delta
  \subseteq \Gamma$. Kita bernalar secara kontrapositif: andaikan $!A
  \notin \Delta$ untuk menunjukkan $\Box!A \notin \Gamma$. Jika
  $!A \notin \Delta$, maka berdasarkan kemaksimalan $\lnot!A \in \Delta$,
  sehingga berdasarkan hipotesis $\Diamond\lnot!A \in \Gamma$. Namun, dalam
  suatu logika modal normal, $\Diamond\lnot!A$ ekuivalen dengan
  $\lnot\Box !A$; dan jika yang disebut terakhir berada dalam $\Gamma$,
  berdasarkan konsistensi $\Box!A \notin\Gamma$, seperti yang diperlukan.
\end{proof}

\iftag{prvBox}{%
  \iftag{prvDiamond}{%
% Box and Diamond are both primitive; prove
% prop:diamond using lem:box-iff-diamond

\begin{prop}\ollabel{prop:diamond}
  Jika $\Gamma$ lengkap dan $\Sigma$-konsisten, maka $\Diamond !A \in
  \Gamma$ jika dan hanya jika terdapat suatu himpunan lengkap dan
  $\Sigma$-konsisten~$\Delta$ sedemikian sehingga $\Diamond\Delta \subseteq
  \Gamma$ dan $!A \in \Delta$.
\end{prop}

\begin{proof}
Andaikan $\Gamma$ lengkap dan $\Sigma$-konsisten. $\Diamond!A \in
\Gamma$ jika dan hanya jika $\lnot\Box\lnot !A \in \Gamma$ berdasarkan
\Dual{} dan ketertutupan. $\lnot\Box\lnot !A \in \Gamma$ jika dan hanya jika
$\Box\lnot !A \notin \Gamma$ berdasarkan
\olref[ccs]{prop:ccs-properties}\olref[ccs]{prop:ccs-lnot}, sebab
$\Gamma$ lengkap dan $\Sigma$-konsisten. Berdasarkan \olref{prop:box},
$\Box\lnot !A \notin \Gamma$ jika dan hanya jika terdapat suatu himpunan
lengkap dan $\Sigma$-konsisten~$\Delta$ dengan
$\Box^{-1}\Gamma \subseteq \Delta$ dan $\lnot!A \notin \Delta$. Sekarang
tinjau sebarang~$\Delta$ semacam itu.
Berdasarkan \olref{lem:box-iff-diamond}, $\Box^{-1}\Gamma \subseteq \Delta$
jika dan hanya jika $\Diamond\Delta \subseteq \Gamma$. Selain itu,
$\lnot !A \notin \Delta$ jika dan hanya jika $!A \in \Delta$ berdasarkan
\olref[ccs]{prop:ccs-properties}\olref[ccs]{prop:ccs-lnot}. Jadi,
$\Diamond!A \in \Gamma$ jika dan hanya jika terdapat suatu himpunan lengkap
dan $\Sigma$-konsisten~$\Delta$ dengan $\Diamond\Delta \subseteq \Gamma$
dan $!A \in \Delta$.
\end{proof}

}{}}{}

\begin{prob}
  Tunjukkan bahwa jika $\Gamma$ lengkap dan $\Sigma$-konsisten, maka
  \iftag{prvBox}{$\Diamond !A \in \Gamma$ jika dan hanya jika terdapat suatu
    himpunan lengkap dan $\Sigma$-konsisten~$\Delta$ sedemikian sehingga
    $\Box^{-1}\Gamma \subseteq \Delta$ dan $!A \in \Delta$.}{$\Box !A \in
    \Gamma$ jika dan hanya jika, bagi setiap himpunan lengkap dan
    $\Sigma$-konsisten~$\Delta$ sedemikian sehingga $\Box^{-1} \Gamma
    \subseteq \Delta$, berlaku $!A \in \Delta$.} \emph{Kerjakan tanpa
    menggunakan \olref[nml][com][mod]{lem:box-iff-diamond}}.
\end{prob}

\end{document}
